\scanpage{60}
\begin{center}atkin page 4\mspage{4}\end{center}

If we now consider the possibility of extending some form of Theorem 1 to $\handwritten{\underline}{C}$(rather than $\handwritten{\underline}{C}^{0}$ )then there is no problem with (H1) and (H2). However\deleted{\unsure{,}} (H4)\handwritten{,}which is essential\handwritten{,}breaks down since the only finite-dimensional subspaces of $\handwritten{\underline}{C}$ which are mapped into themselves by all the $\handwritten{\underline}{T_p}$ are subspaces of $\handwritten{\underline}{C}^{\handwritten{+}}$. Indeed any modular form with a pole\handwritten{\tikz[baseline,overlay,x=1em,y=1em]{\draw[very thin,gray] (-.9,-.35) .. controls (-.7,-.63) and (.85,-.8) .. (1.3,-1) .. controls (1.55,-1.1) and (1.39,-.59) .. (1.58,-.66) .. controls (1.76,-.79) and (1.63,-.92) .. (1.48,-1.06);}} gives rise to an infinite-dimensional subspace by the repeated action of only one $\handwritten{\underline}{T_p}$ .

The essential distinction,then,lies between $\handwritten{\underline}{C}$ and $\handwritten{\underline}{C}^{\handwritten{0}}$. It so happens that $\handwritten{\underline}{C}^{\handwritten{0}}$ is null for $\handwritten{\underline}{k}\handwritten{\geqslant}0$ and there is thus an association between Theorem 1 and forms of negative dimension. In fact such theory as we are able to develop for $\handwritten{\underline}{C}$ is viable for forms of any dimension,but the author has mainly pursued it for positive and zero dimension where nothing else is available.

We now choose the simplest of all modular forms,in fact the modular function($\handwritten{\underline}{k}=0$) $\handwritten{\underline}{j(z)}$ which is the Hauptmodul of $\handwritten{\underline}{\handwritten{\Gamma}}$ defined by
\[
\begin{aligned}
j(z)&=\handwritten{\left\{1+240\sum_{n=1}^{\infty}\underline{\sigma}_{3}(n)\,x^n\right\}^{3}\Bigm/ x\prod_{r=1}^{\infty}(1-x^{r})^{24}\unsure{.}}\\
&=x^{-1}+744+\handwritten{\sum_{n=1}^{\infty}}c(n)\,x^n,
\end{aligned}
\]
\reversemarginpar\margin{\handwritten{Sigma}}\normalmarginpar
and consider what kind of property
\[
a(1)=1,\ a(np)-a(n)a(p)+p^{-1}a(n/p)=0\quad (\text{all prime }p,\text{all integral }n),
\tag*{(2)}
\]
can exist for $\handwritten{\underline}{a(n)}$ associated in some way with $\handwritten{\underline}{c(n)}$. Nothing identical is possible,for the reasons already given. However (2) makes sense if the $\handwritten{\underline}{a(n)}$ are considered not\inserted{necessarily} as complex numbers but as elements of any commutative ring with identity; since in fact the $\handwritten{\underline}{a(n)}$ of Theorem 1 are algebraic integers(as shown by Rankin and Rushforth\handwritten{\uwave}{(9)}) a natural choice of ring is $\handwritten{\underline}{Z/q^{\handwritten{\underline}{\alpha}}Z}$\margin{\handwritten{alpha}} ,that is, the integers modulo $\handwritten{\underline}{q}^{\handwritten{\underline}{\alpha}}$,for some fixed prime $\handwritten{\underline}{q}$. We cannot expect some \handwritten{\uline}{fixed} $\handwritten{\underline}{a(n)}$ to be involved for all $\handwritten{\underline}{q}$,or even for fixed $\handwritten{\underline}{q}$ and all $\handwritten{\underline}{\alpha}$,for then (2) would be identical again. Thus we may expect $\handwritten{\underline}{a(n)}$ to depend on $\handwritten{\underline}{q}$ and $\handwritten{\underline}{\alpha}$ . At this point we \deleted{\illegible{1}} state some actual results,which indicate that something of this kind actually happens; for $\handwritten{\underline}{q}=2,3,5,7$,and $13$ we can prove for all $\handwritten{\underline}{\alpha}$,and otherwise.for a few fixed $a$ the following

% archive sheet number: 59
